\section{봉인된 다중 gate 설계}
\label{sec:design}

\subsection{Cohort와 one-shot access}

Table~\ref{tab:gates}는 세 confirmatory cohort를 요약한다. 각 gate에는
development generator, 별도의 encrypted seed commitment, zero-access audit,
frozen source digest와 atomic one-shot lock이 있다. Sealed result를 threshold,
feature, model population 또는 sample size 변경에 재사용하지 않는다.

\begin{table}[ht]
\centering
\caption{Confirmatory cohort. Run 수는 world 수, 여섯 control과 세 optimizer
seed의 곱이다.}
\label{tab:gates}
\begin{tabular}{lrrrl}
\toprule
Gate & Worlds & Runs & 최대 horizon & Primary 목적 \\
\midrule
Structural OOD & 50 & 900 & 1 & unseen mechanism과 routing \\
Open-loop rollout & 62 & 1,116 & 32 & recursive structural error \\
Criterion validity & 69 & 1,242 & 8 & plan failure와 first failure \\
\bottomrule
\end{tabular}
\end{table}

OOD와 rollout gate는 paired comparison에 여섯 control 전체를 사용한다.
Criterion-validity gate는 nondegenerate error를 생성하는 네 control에서 hazard
model을 fit한다. Dense와 oracle-signature control은 mandatory structural-zero
check로 유지하며 report에서 제거하지 않는다.

\subsection{Structural OOD gate}

Development mechanism과 sealed mechanism은 world level에서 disjoint하다.
Generator는 unseen recombination, unseen structural mode, unseen mechanism
parameter, unseen action composition, bridge-consistent shift와
bridge-violating negative control을 구분한다. Primary comparison은 capacity와
compute를 유지하면서 correct routing을 random/wrong routing과 pair한다.
Frozen power rule로 sealed world 50개를 선택했다.

\subsection{Open-loop gate}

Rollout protocol은 horizon \(1,2,4,8,16,32\)에서 teacher-forced prediction과
recursive prediction을 분리한다. \(\DeltaTheta\) trajectory의 area,
fixed-layer vector, structural failure 없는 survival, bridge violation과
tracking diagnostic을 기록한다. 누적 worst-case numerical bound가
vacuous하더라도 recursive inequality audit를 유지한다.

\subsection{Criterion-validity gate}

각 world는 32개 unit을 가진다. 한 unit에는 domain-separated 8-step probe와
\[
(1,2,4,8,1,2,4,8)
\]
horizon을 가진 task 8개가 있다. Task utility는 structural predicate만으로
생성한다. Primary population은 layer-routed, strict-factorized,
random-routed, wrong-routed control이다.

Scalar baseline은 training negative log likelihood, one-step/open-loop latent
mean-squared error, endpoint exactness, task horizon과 goal weight, model
identity, plan-selection margin, selected-versus-alternative scalar error
contrast를 포함한다. Augmented model은 prespecified task-local
\(\DeltaTheta\), fixed-total, fixed-layer와 goal-aligned structural contrast를
추가한다. Mandatory non-gating sensitivity는 baseline에 direct raw layer
mismatch를 제공한다.

Ridge logistic discrete hazard \(h_i(t)\)를 development world에서 fit한다.
Cumulative failure probability는
\[
\widehat F_i(t)
=1-\prod_{s=1}^{t}\{1-\widehat h_i(s)\}
\]
이다. 모든 coefficient, feature order, mean, scale과 penalty를 sealed seed
reveal 전에 동결한다.

\subsection{Estimand와 inference}

두 co-primary effect는
\[
\begin{split}
\mathcal E_{\mathrm{bin}}
&=
\operatorname{BS}_{\mathrm{scalar}}(8)
-
\operatorname{BS}_{\mathrm{TSI}}(8),\\
\mathcal E_{\mathrm{int}}
&=
\frac{1}{8}\sum_{t=1}^{8}
\left[
\operatorname{BS}_{\mathrm{scalar}}(t)
-
\operatorname{BS}_{\mathrm{TSI}}(t)
\right].
\end{split}
\]
이다. 여기서 \(\operatorname{BS}(8)\)의 8은 task identifier가 아니라 고정된
8-step candidate trajectory의 terminal time index \(t=8\)이다. Integrated
estimand는 이 trajectory의 8개 time index를 평균한다. 양수이면 structural
augmentation이 우세하다. Unit-level score는
world 안에서 primary control과 nested optimizer seed에 대해 먼저 평균한다.
One-sided world-level Student test에는 Holm familywise correction을 적용한다
\citep{holm1979simple}. 또한 20,000-replicate world-cluster bootstrap과
world-random-intercept, seed-nested variance decomposition에서 Bonferroni
simultaneous lower bound가 양수여야 한다. 두 point effect는 prespecified
smallest effect size of interest \(\SESOI=0.0025\)를 만족해야 한다. Inferential
null은 nonpositive effect이며 confidence lower bound가 \(\SESOI\)를 넘는다고
주장하지 않는다.
